\documentclass[11pt]{article}

\usepackage[T1]{fontenc}
\usepackage{lmodern}
\usepackage[margin=1in]{geometry}
\usepackage{amsmath,amssymb,amsthm}
\usepackage{booktabs,longtable}
\usepackage[hidelinks]{hyperref}

\newtheorem{theorem}{Theorem}
\newtheorem{lemma}[theorem]{Lemma}
\theoremstyle{definition}
\newtheorem{definition}[theorem]{Definition}
\theoremstyle{remark}
\newtheorem{remark}[theorem]{Remark}

\newcommand{\F}{\mathbb{F}}
\newcommand{\Sp}{\mathcal{S}}
\newcommand{\Aut}{\operatorname{Aut}}

\title{Improving the covering code bound $K(1,7,1)\le 48$\\ requires an asymmetric code}
\author{Noah Parsons\\ \small Independent researcher\thanks{%
\texttt{parsons.m.noah@gmail.com}, ORCID \href{https://orcid.org/0009-0000-7224-6040}{0009-0000-7224-6040}.}}
\date{October 2026}

\begin{document}
\maketitle

\begin{abstract}
Let $K(1,7,1)$ be the smallest number of words in a covering code of radius $1$ in the mixed
Hamming space $\F_3\times\F_2^7$. The tables of mixed binary/ternary covering codes give
$42\le K(1,7,1)\le 48$, and the upper bound comes from the direct product of $\F_3$ with the
binary Hamming code of length $7$. We prove, with the help of a computer, that every covering
code of radius $1$ in this space with at most $47$ words has trivial automorphism group.
The proof reduces the question to one representative of each of the $46$ conjugacy classes of
isometries of prime order and solves an integer program for each; the minimum size of an
invariant covering code turns out to be $48$, $54$ or $56$ in every case. Two independent
integer-programming solvers agree on all $46$ cases.
\end{abstract}

\section{Introduction}

A \emph{covering code} of radius $R$ in a Hamming space is a set of words such that every word
of the space is within Hamming distance $R$ of at least one of them. For mixed spaces
$\F_3^t\times\F_2^b$ the smallest possible size is denoted $K(t,b,R)$; bounds are collected in
tables maintained by K\'eri~\cite{keri}, following
\cite{ostergard-hamalainen-1997,bertolo-ostergard-weakley-2004}, and the general theory is in the
monograph~\cite{cohen-book}.

This note concerns the entry $K(1,7,1)$, for which the tables give
\begin{equation}\label{eq:bounds}
42 \le K(1,7,1) \le 48 .
\end{equation}
The lower bound is due to Lang, Quistorff and Schneider~\cite{lqs-2007}. The upper bound is the
elementary construction $K(t,b,R)\le 3K(t-1,b,R)$: if $H\subseteq\F_2^7$ is the Hamming code,
which has $16$ words and covering radius $1$, then $\F_3\times H$ has $48$ words and covering
radius $1$. This code is highly symmetric.

A standard way to look for good codes is to prescribe a group of symmetries and search only
among codes it leaves invariant; Östergård and Weakley used this to construct covering
codes~\cite{ostergard-weakley-1999}. Here the same idea is used in the opposite direction: to
show that a code improving the upper bound in~\eqref{eq:bounds} cannot have any symmetry at all.

\begin{theorem}\label{thm:main}
Let $C\subseteq\F_3\times\F_2^7$ be a covering code of radius $1$ with $|C|\le 47$. Then the
only isometry of $\F_3\times\F_2^7$ that maps $C$ onto itself is the identity.
\end{theorem}

The proof is computer-assisted (Section~\ref{sec:computation}). It also determines, for each
type of symmetry of prime order, the smallest covering code having that symmetry
(Table~\ref{tab:classes}).

% TODO(Noah): before submitting, read [ostergard-hamalainen-1997] and
% [bertolo-ostergard-weakley-2004] in full and check that neither already states Theorem 1.

\section{Preliminaries}

Write $\Sp=\F_3\times\F_2^7$; it has $3\cdot 2^7=384$ words. The Hamming distance $d(x,y)$ is the
number of coordinates in which $x$ and $y$ differ, and $B(x)=\{y\in\Sp: d(x,y)\le 1\}$. Each
ball has $1+2+7=10$ words.

\begin{definition}
An \emph{isometry} of $\Sp$ is a bijection $\Sp\to\Sp$ that preserves $d$. The isometries form a
group $G$. For a code $C\subseteq\Sp$, its \emph{automorphism group} is
$\Aut(C)=\{g\in G: g(C)=C\}$.
\end{definition}

The isometries of a Hamming space are exactly the maps that permute coordinates carrying
alphabets of the same size and apply a permutation of the alphabet in each
coordinate~\cite{cohen-book}. Here the ternary coordinate cannot be moved, so
\[
G \;\cong\; S_3 \times \bigl(S_2 \wr S_7\bigr),
\qquad |G| = 6\cdot 2^7\cdot 7! = 3\,870\,720 .
\]
An element of $S_2\wr S_7$ is a pair $(f,\sigma)$: a permutation $\sigma$ of the seven binary
coordinates followed by flipping the coordinates in a set $f$.

\section{Reduction to finitely many integer programs}

\begin{lemma}\label{lem:conj}
If $C$ is a covering code of radius $1$ and $h\in G$, then $h(C)$ is a covering code of radius
$1$ with $|h(C)|=|C|$, and $\Aut(h(C)) = h\,\Aut(C)\,h^{-1}$.
\end{lemma}

\begin{proof}
Since $h$ is a bijection preserving distance, $h(B(c))=B(h(c))$, so the balls around $h(C)$ are
the images of the balls around $C$ and still cover $\Sp$. Further, $g(C)=C$ if and only if
$(hgh^{-1})(h(C))=h(C)$.
\end{proof}

\begin{lemma}\label{lem:prime}
To prove Theorem~\ref{thm:main} it suffices to show: for one representative $r$ of each
conjugacy class of elements of prime order in $G$, no covering code $C$ of radius $1$ with
$|C|\le 47$ satisfies $r(C)=C$.
\end{lemma}

\begin{proof}
Suppose $|C|\le 47$ and $\Aut(C)$ contains some $g_0\neq\mathrm{id}$. Since $G$ is finite, $g_0$
has some order $n>1$. Choose a prime $p$ dividing $n$; then $g=g_0^{\,n/p}$ has order exactly
$p$, and $g\in\Aut(C)$ because every power of an automorphism is an automorphism.
Let $r=hgh^{-1}$ be the chosen representative
of the conjugacy class of $g$. By Lemma~\ref{lem:conj}, $h(C)$ is a covering code of the same
size and $r\in\Aut(h(C))$, contradicting the assumption.
\end{proof}

\begin{lemma}\label{lem:classes}
The elements of prime order in $G$ fall into exactly $46$ conjugacy classes: $39$ of order $2$,
$5$ of order $3$, one of order $5$ and one of order $7$. Representatives are $(\alpha,(f,\sigma))$
where
\begin{itemize}
\item $p=2$: $\alpha$ is the identity or a transposition of $\F_3$; $\sigma$ is a product of
  $c\in\{0,1,2,3\}$ disjoint transpositions; $f$ flips $k$ of the $7-2c$ coordinates fixed by
  $\sigma$, where $0\le k\le 7-2c$; and not all of $\alpha,c,k$ are trivial;
\item $p=3$: $\alpha$ is the identity or the $3$-cycle $x\mapsto x+1$; $\sigma$ is a product of
  $c\in\{0,1,2\}$ disjoint $3$-cycles; $f=\emptyset$; and not both $\alpha$ and $c$ are trivial;
\item $p=5$: $\alpha=\mathrm{id}$, $\sigma$ a $5$-cycle, $f=\emptyset$;
\item $p=7$: $\alpha=\mathrm{id}$, $\sigma$ a $7$-cycle, $f=\emptyset$.
\end{itemize}
\end{lemma}

\begin{proof}
An element $(\alpha,\beta)$ of the direct product has prime order $p$ exactly when
$\alpha^p=1$, $\beta^p=1$ and they are not both the identity, and it is conjugate to
$(\alpha',\beta')$ exactly when $\alpha$ is conjugate to $\alpha'$ in $S_3$ and $\beta$ to
$\beta'$ in $S_2\wr S_7$. In $S_3$ the classes are the identity, the transpositions and the
$3$-cycles.

Conjugacy classes in $S_2\wr S_7$ are determined by the cycle type of $\sigma$ together with,
for each cycle, the parity of the number of flipped coordinates on it~\cite[Ch.~4]{james-kerber}.
If $\beta=(f,\sigma)$ has $\beta^p=1$ then $\sigma^p=1$, so every cycle of $\sigma$ has length $1$
or $p$. On a cycle of length $p$, $\beta^p$ flips each coordinate of the cycle once for every
flip on the cycle, so the parity must be even; such a cycle is conjugate to one with no flips. On
a fixed coordinate a flip has order $2$, so flipped fixed coordinates occur only when $p=2$.
This gives the representatives listed, and counting them gives
$2\cdot(8+6+4+2)-1=39$ classes of order $2$, $3+2=5$ of order $3$, and one each of orders $5$
and $7$. As an independent check, the program computes the same counts by brute-force
enumeration of the whole group on five smaller spaces (\texttt{tests/test\_classify.py}).
\end{proof}

\paragraph{The integer program.} Fix a representative $r$, and let $O_1,\dots,O_m$ be the orbits
of the cyclic group $\langle r\rangle$ on $\Sp$, with representatives $w_i\in O_i$. A code is
$r$-invariant exactly when it is a union of orbits. If $C$ is $r$-invariant then a word $w$ is
covered if and only if $r(w)$ is, because $r$ is an isometry; so it suffices to check coverage
at $w_1,\dots,w_m$. Hence the smallest $r$-invariant covering code has size
\begin{equation}\label{eq:ilp}
m(r) \;=\; \min\Bigl\{ \sum_{j=1}^m |O_j|\,y_j \;:\;
\sum_{j:\,O_j\cap B(w_i)\neq\emptyset} y_j \ge 1 \ (1\le i\le m),\ y\in\{0,1\}^m \Bigr\}.
\end{equation}
By Lemmas~\ref{lem:prime} and~\ref{lem:classes}, Theorem~\ref{thm:main} is equivalent to
$m(r)\ge 48$ for all $46$ representatives.

\section{Computation}\label{sec:computation}

The programs~\eqref{eq:ilp} have between $60$ and $288$ variables. They were solved twice, by
separately written code:
\begin{enumerate}
\item \textbf{First pass.} \texttt{src/classify.py} builds each program from the class
  representative and solves it with HiGHS (through SciPy~1.16.3), either minimising and reading
  the proven lower bound or adding the constraint $\sum|O_j|y_j\le 47$ and obtaining
  infeasibility. All $46$ classes were excluded; total solver time $2655$\,s.
\item \textbf{Second pass.} \texttt{src/recheck.py}, which shares no code with the first pass,
  replaces each representative $r$ by a random conjugate $hrh^{-1}$ (seed $2026$), rebuilds the
  orbits and the program from scratch, and minimises~\eqref{eq:ilp} with SCIP~10.0
  through PySCIPOpt~6.3.0. SCIP proved an optimal value $m(r)\in\{48,54,56\}$ in all $46$ cases;
  total solver time $1715$\,s on one core of an Intel Core Ultra~7 255HX.
\end{enumerate}
Before the second pass, the method was tested on a case with a known answer: for two classes
of order $2$, allowing $48$ words, SCIP produced an invariant covering code of size $48$, and an
independent brute-force checker (\texttt{src/checker.py}) confirms that both codes cover all
$384$ words (\texttt{results/codes/K\_1\_7\_1\_M48\_scip\_*.txt}). Partial runs with two
further solvers, CBC on $14$ classes and HiGHS on a random conjugate for $11$, agree as well.
Table~\ref{tab:classes} lists every class with both results.

\begin{remark}
Both solvers work in floating-point arithmetic, so their verdicts are strong evidence rather
than certificates of the kind an exact verifier would produce. Two solvers written by different
teams, applied to independently built instances, agree on every case.
\end{remark}

\section{Discussion}

For many classes $m(r)=48$: the record $48$ can be reached by codes with many different
symmetries, but not one of them can be shrunk while keeping its symmetry. Theorem~\ref{thm:main}
therefore says that a search for a $47$-word code can ignore all symmetric codes, which is where
searches usually start. Whether $K(1,7,1)=48$ remains open; deciding it requires handling the
asymmetric codes, for which the reduction above gives no help.

\paragraph{Code and data.} All programs, logs and the generator of Table~\ref{tab:classes} are at
\url{https://github.com/Noah-Parsons/covering-codes}.

\begin{thebibliography}{9}

\bibitem{bertolo-ostergard-weakley-2004}
R.~Bertolo, P.~R.~J. Östergård and W.~D. Weakley,
An updated table of binary/ternary mixed covering codes,
\emph{J. Combin. Des.} 12 (2004), 157--176.

\bibitem{cohen-book}
G.~Cohen, I.~Honkala, S.~Litsyn and A.~Lobstein,
\emph{Covering Codes}, North-Holland, Amsterdam, 1997.

\bibitem{james-kerber}
G.~James and A.~Kerber,
\emph{The Representation Theory of the Symmetric Group},
Addison-Wesley, Reading, MA, 1981.

\bibitem{keri}
G.~Kéri, Tables for bounds on covering codes,
\url{https://old.sztaki.hu/~keri/codes/index.htm} (accessed 29 September 2026).

\bibitem{lqs-2007}
W.~Lang, J.~Quistorff and E.~Schneider,
New results on integer programming for codes,
\emph{Congr. Numer.} 188 (2007), 97--107.

\bibitem{ostergard-hamalainen-1997}
P.~R.~J. Östergård and H.~O. Hämäläinen,
A new table of binary/ternary mixed covering codes,
\emph{Des. Codes Cryptogr.} 11 (1997), 151--178.

\bibitem{ostergard-weakley-1999}
P.~R.~J. Östergård and W.~D. Weakley,
Constructing covering codes with given automorphisms,
\emph{Des. Codes Cryptogr.} 16 (1999), 65--73.

\end{thebibliography}

\clearpage
\appendix
\section{The 46 classes}

\begin{center}
\small
\begin{longtable}{rrcrrrrrrr}
\caption{One row per conjugacy class of prime-order isometries. $\alpha$: map on the ternary
coordinate; $c$: number of $p$-cycles on binary coordinates; $k$: flipped fixed binary
coordinates; orbits: number of variables in~\eqref{eq:ilp}. HiGHS: proven lower bound when the
first pass minimised (``--'' where it proved the capped program infeasible). SCIP: proven
optimum $m(r)$ on a random conjugate. Times in seconds.}\label{tab:classes}\\
\toprule
\# & $p$ & $\alpha$ & $c$ & $k$ & orbits & HiGHS & time & SCIP $m(r)$ & time \\
\midrule
\endfirsthead
\toprule
\# & $p$ & $\alpha$ & $c$ & $k$ & orbits & HiGHS & time & SCIP $m(r)$ & time \\
\midrule
\endhead
% generated by paper/make_table.py -- do not edit by hand
1 & 7 & id & 1 & 0 & 60 & -- & 0.0 & 48 & 0.1 \\
2 & 5 & id & 1 & 0 & 96 & -- & 0.0 & 56 & 0.4 \\
3 & 3 & $x\mapsto x+1$ & 0 & 0 & 128 & -- & 0.0 & 48 & 0.0 \\
4 & 3 & $x\mapsto x+1$ & 1 & 0 & 128 & -- & 0.7 & 48 & 1.3 \\
5 & 3 & $x\mapsto x+1$ & 2 & 0 & 128 & -- & 1.2 & 48 & 1.6 \\
6 & 3 & id & 1 & 0 & 192 & -- & 3.7 & 54 & 41.9 \\
7 & 3 & id & 2 & 0 & 144 & -- & 5.6 & 48 & 1.8 \\
8 & 2 & $(1\,2)$ & 0 & 0 & 256 & 48.00 & 11.0 & 48 & 6.1 \\
9 & 2 & $(1\,2)$ & 0 & 1 & 192 & 48.00 & 3.1 & 48 & 2.3 \\
10 & 2 & $(1\,2)$ & 0 & 2 & 192 & 48.00 & 17.6 & 48 & 10.4 \\
11 & 2 & $(1\,2)$ & 0 & 3 & 192 & 48.00 & 163.7 & 48 & 24.9 \\
12 & 2 & $(1\,2)$ & 0 & 4 & 192 & 48.00 & 48.4 & 48 & 19.0 \\
13 & 2 & $(1\,2)$ & 0 & 5 & 192 & 48.00 & 92.6 & 48 & 25.1 \\
14 & 2 & $(1\,2)$ & 0 & 6 & 192 & 48.00 & 56.2 & 48 & 24.4 \\
15 & 2 & $(1\,2)$ & 0 & 7 & 192 & 48.00 & 97.0 & 48 & 26.1 \\
16 & 2 & $(1\,2)$ & 1 & 0 & 224 & 48.00 & 41.4 & 48 & 30.0 \\
17 & 2 & $(1\,2)$ & 1 & 1 & 192 & 48.00 & 17.5 & 48 & 13.1 \\
18 & 2 & $(1\,2)$ & 1 & 2 & 192 & 48.00 & 83.0 & 48 & 28.7 \\
19 & 2 & $(1\,2)$ & 1 & 3 & 192 & 48.00 & 93.2 & 48 & 35.4 \\
20 & 2 & $(1\,2)$ & 1 & 4 & 192 & 48.00 & 73.5 & 48 & 70.0 \\
21 & 2 & $(1\,2)$ & 1 & 5 & 192 & 48.00 & 100.7 & 48 & 30.9 \\
22 & 2 & $(1\,2)$ & 2 & 0 & 208 & 48.00 & 173.7 & 48 & 21.9 \\
23 & 2 & $(1\,2)$ & 2 & 1 & 192 & 48.00 & 43.0 & 48 & 14.2 \\
24 & 2 & $(1\,2)$ & 2 & 2 & 192 & 48.00 & 107.3 & 48 & 35.8 \\
25 & 2 & $(1\,2)$ & 2 & 3 & 192 & 48.00 & 133.5 & 48 & 33.8 \\
26 & 2 & $(1\,2)$ & 3 & 0 & 200 & 48.00 & 42.5 & 48 & 52.5 \\
27 & 2 & $(1\,2)$ & 3 & 1 & 192 & 48.00 & 22.6 & 48 & 35.8 \\
28 & 2 & id & 0 & 1 & 192 & -- & 1.7 & 48 & 2.0 \\
29 & 2 & id & 0 & 2 & 192 & -- & 1.7 & 48 & 2.2 \\
30 & 2 & id & 0 & 3 & 192 & 48.00 & 37.3 & 48 & 23.0 \\
31 & 2 & id & 0 & 4 & 192 & 48.00 & 29.8 & 48 & 14.7 \\
32 & 2 & id & 0 & 5 & 192 & 48.00 & 84.8 & 48 & 37.0 \\
33 & 2 & id & 0 & 6 & 192 & 48.00 & 53.8 & 48 & 26.6 \\
34 & 2 & id & 0 & 7 & 192 & 48.00 & 56.0 & 48 & 11.2 \\
35 & 2 & id & 1 & 0 & 288 & 48.00 & 192.3 & 48 & 626.4 \\
36 & 2 & id & 1 & 1 & 192 & -- & 6.2 & 48 & 16.0 \\
37 & 2 & id & 1 & 2 & 192 & -- & 7.6 & 48 & 16.4 \\
38 & 2 & id & 1 & 3 & 192 & 48.00 & 62.3 & 48 & 52.5 \\
39 & 2 & id & 1 & 4 & 192 & 48.00 & 94.9 & 48 & 24.1 \\
40 & 2 & id & 1 & 5 & 192 & 48.00 & 75.7 & 48 & 60.5 \\
41 & 2 & id & 2 & 0 & 240 & 48.00 & 176.6 & 48 & 70.9 \\
42 & 2 & id & 2 & 1 & 192 & 48.00 & 34.9 & 48 & 15.4 \\
43 & 2 & id & 2 & 2 & 192 & 48.00 & 25.2 & 48 & 14.7 \\
44 & 2 & id & 2 & 3 & 192 & 48.00 & 119.7 & 48 & 40.5 \\
45 & 2 & id & 3 & 0 & 216 & 48.00 & 105.2 & 48 & 27.7 \\
46 & 2 & id & 3 & 1 & 192 & 48.00 & 57.0 & 48 & 45.8 \\
\bottomrule

\end{longtable}
\end{center}

\end{document}
